\documentclass[a4paper]{article}

\usepackage[spanish]{babel}
\usepackage{graphicx}
\usepackage{amsmath, amssymb, mathrsfs}
\usepackage[margin=2cm]{geometry}
\usepackage{fancyhdr}
\usepackage{enumerate}
\usepackage[shortlabels]{enumitem}
\usepackage{parskip}
\usepackage[most]{tcolorbox}
\usepackage[hidelinks]{hyperref}
\usepackage{float}
\usepackage{booktabs}
\usepackage{pgfplots}
\pgfplotsset{compat=1.18}
\usepgfplotslibrary{groupplots}
\usetikzlibrary{arrows.meta}



% cabecera
\pagestyle{fancy}
\fancyhead[l]{Física Matemática I}
\fancyhead[c]{Taller 10}

\setlength{\headheight}{14.5pt}
\renewcommand{\headrulewidth}{0.2pt} % linea horizontal

\begin{document}
\pagenumbering{roman}

\begin{titlepage}
  \centering
  {\Large Universidad de Antioquia\par}
  \vspace{2cm}
  {\huge\bfseries Física Matemática I\par}
  \vspace{0.4cm}
  {\Large Solución al Taller 10 \par}
  \vspace{1.2cm}
  \begin{tcolorbox}[
    width=0.86\textwidth,
    colback=white,
    colframe=black!55,
    arc=2mm,
    boxrule=0.5pt,
    title={\centering Algunas ecuaciones del taller},
    fonttitle=\bfseries
  ]
  \begin{align*}
    g(x,t) &= e^{2xt-t^2}=\sum_{n=0}^{\infty}\frac{H_n(x)}{n!}t^n, \\
    H_n(x) &= (-1)^n e^{x^2}\frac{d^n}{dx^n}\left(e^{-x^2}\right), \\
    H'_n(x) &= 2nH_{n-1}(x),
    \qquad
    2xH_n(x)=2nH_{n-1}(x)+H_{n+1}(x), \\
    \int_{-\infty}^{\infty} e^{-x^2}H_m(x)H_n(x)\,dx
    &= \sqrt{\pi}\,2^n n!\,\delta_{mn}, \\
    H_n(-x) &= (-1)^n H_n(x), \\
    H''_n(x)-2xH'_n(x)+2nH_n(x) &= 0, \\
    x^{2r} &= \frac{(2r)!}{2^{2r}}
    \sum_{k=0}^{r}\frac{H_{2k}(x)}{(2k)!\,(r-k)!}, \\
    \psi_n(x) &= \left(\frac{\alpha}{\sqrt{\pi}\,2^n n!}\right)^{1/2}
    H_n(\alpha x)e^{-\alpha^2x^2/2}, \\
    E_n &= \hbar\omega\left(n+\frac{1}{2}\right).
  \end{align*}
  \end{tcolorbox}
  \vspace{1.2cm}
  {\large Autor:\par}
  \vspace{0.2cm}
  {\large Daniel Soto Villada\par}
  {\large Estudiante de Astronomía\par}
  \vfill
  {\large \today\par}
\end{titlepage}
\setcounter{page}{2}

\newpage
\tableofcontents
\newpage
\pagenumbering{arabic}


\section{Problema 1}
Utilizando la función generatriz de los polinomios de Hermite
$$
g(x, t) = e^{2xt - t^2} = \sum_{n=0}^{\infty} \frac{H_n(x)}{n!} t^n
$$
muestre las siguientes relaciones (para ello derive la función generatriz con respecto a $x$ o $t$):
\begin{itemize}
    \item[a)] $2xH_n(x) = 2nH_{n-1}(x) + H_{n+1}(x)$.
    \item[b)] $H'_n(x) = 2nH_{n-1}(x)$.
    \item[c)] $H_{2n+1}(0) = 0$, $H_{2n}(0) = (-1)^n \frac{(2n)!}{n!}$.
    \item[d)] Usando a) y b), muestre que $H''_n(x) - 2xH'_n(x) + 2nH_n(x) = 0$.
    \item[e)] $e^{2x} = e \sum_{n=0}^{\infty} \frac{H_n(x)}{n!}$.
\end{itemize}

\subsection{Solución}
Para resolver este problema nos basamos en la teoría de los polinomios de Hermite presentada en el Capítulo 8, Sección 8.5 del texto de Alonso Sepúlveda. La función generatriz de los polinomios de Hermite se define como (Alonso Sepúlveda, Sección 8.5, p. 338):
$$
g(x, t) = e^{2xt - t^2} = \sum_{n=0}^{\infty} \frac{H_n(x)}{n!} t^n, \quad \text{para } |t| < \infty.
$$
A partir de esta expresión fundamental, todas las relaciones de recurrencia, derivadas y valores específicos solicitados pueden obtenerse mediante diferenciación parcial respecto a $x$ y $t$, seguida de la igualación de coeficientes de potencias iguales de $t$. A continuación, se detalla el desarrollo analítico paso a paso.

Antes de entrar en las demostraciones algebraicas, conviene visualizar dos aspectos importantes: la forma de los primeros polinomios de Hermite y el comportamiento de la función generatriz para distintos valores del parámetro auxiliar $t$.

\begin{center}
\begin{tikzpicture}
\begin{axis}[
    width=0.78\textwidth,
    height=6cm,
    xmin=-3, xmax=3,
    ymin=-1.15, ymax=1.15,
    axis lines=middle,
    grid=both,
    grid style={black!12},
    xlabel={$x$},
    ylabel={valor normalizado},
    title={Primeros polinomios de Hermite},
    legend style={font=\scriptsize, cells={anchor=west}, at={(0.02,0.98)}, anchor=north west},
    samples=220,
    domain=-3:3
]
\addplot[blue, thick] {1};
\addlegendentry{$H_0(x)$}
\addplot[red, thick] {x/3};
\addlegendentry{$H_1(x)/6$}
\addplot[green!55!black, thick] {(4*x^2-2)/34};
\addlegendentry{$H_2(x)/34$}
\addplot[orange, thick] {(8*x^3-12*x)/180};
\addlegendentry{$H_3(x)/180$}
\addplot[purple, thick] {(16*x^4-48*x^2+12)/876};
\addlegendentry{$H_4(x)/876$}
\end{axis}
\end{tikzpicture}

\small\textbf{Figura 1.} Primeros polinomios de Hermite normalizados para comparar su forma y paridad. Los polinomios con $n$ par son funciones pares, mientras que los de $n$ impar son funciones impares.
\end{center}

\begin{center}
\begin{tikzpicture}
\begin{axis}[
    width=0.78\textwidth,
    height=6cm,
    xmin=-3, xmax=3,
    ymin=0, ymax=20,
    axis lines=left,
    grid=both,
    grid style={black!12},
    xlabel={$x$},
    ylabel={$g(x,t)$},
    title={$g(x,t)=e^{2xt-t^2}$},
    legend style={font=\scriptsize, at={(0.02,0.98)}, anchor=north west},
    samples=220,
    domain=-3:3
]
\addplot[blue, thick] {exp(0.5*x-0.0625)};
\addlegendentry{$t=0.25$}
\addplot[red, thick] {exp(x-0.25)};
\addlegendentry{$t=0.50$}
\addplot[green!55!black, thick] {exp(1.5*x-0.5625)};
\addlegendentry{$t=0.75$}
\addplot[orange, thick] {exp(2*x-1)};
\addlegendentry{$t=1.00$}
\end{axis}
\end{tikzpicture}

\small\textbf{Figura 2.} Función generatriz $g(x,t)=e^{2xt-t^2}$ para varios valores de $t$. Al expandir esta función en potencias de $t$ aparecen como coeficientes los polinomios $H_n(x)$.
\end{center}

\vspace{0.3cm}
\noindent\textbf{Demostración del inciso a):} \\
Derivamos la función generatriz respecto al parámetro $t$:
$$
\frac{\partial g}{\partial t} = \frac{\partial}{\partial t} \left( e^{2xt - t^2} \right) = (2x - 2t) e^{2xt - t^2} = (2x - 2t) g(x, t).
$$
Por otro lado, derivando la serie de potencias término a término respecto a $t$, obtenemos:
$$
\frac{\partial g}{\partial t} = \sum_{n=0}^{\infty} \frac{H_n(x)}{n!} n t^{n-1} = \sum_{n=1}^{\infty} \frac{H_n(x)}{(n-1)!} t^{n-1}.
$$
Realizamos un cambio de índice haciendo $m = n - 1$ (de modo que $n = m + 1$), lo que nos permite reescribir la sumatoria comenzando desde $m = 0$:
$$
\frac{\partial g}{\partial t} = \sum_{m=0}^{\infty} \frac{H_{m+1}(x)}{m!} t^m = \sum_{n=0}^{\infty} \frac{H_{n+1}(x)}{n!} t^n.
$$
Igualamos esta expresión con el resultado de derivar la forma cerrada:
$$
(2x - 2t) \sum_{n=0}^{\infty} \frac{H_n(x)}{n!} t^n = \sum_{n=0}^{\infty} \frac{H_{n+1}(x)}{n!} t^n.
$$
Distribuimos el término $(2x - 2t)$ en el lado izquierdo:
$$
2x \sum_{n=0}^{\infty} \frac{H_n(x)}{n!} t^n - 2 \sum_{n=0}^{\infty} \frac{H_n(x)}{n!} t^{n+1} = \sum_{n=0}^{\infty} \frac{H_{n+1}(x)}{n!} t^n.
$$
En el segundo término de la izquierda, hacemos un cambio de índice $m = n + 1$ (es decir, $n = m - 1$), de modo que la suma comienza en $m = 1$:
$$
2x \sum_{n=0}^{\infty} \frac{H_n(x)}{n!} t^n - 2 \sum_{n=1}^{\infty} \frac{H_{n-1}(x)}{(n-1)!} t^n = \sum_{n=0}^{\infty} \frac{H_{n+1}(x)}{n!} t^n.
$$
Notamos que $\frac{1}{(n-1)!} = \frac{n}{n!}$, por lo que podemos reescribir la segunda sumatoria para que comience en $n=0$ (el término $n=0$ es cero):
$$
\sum_{n=0}^{\infty} \left[ 2x \frac{H_n(x)}{n!} - 2n \frac{H_{n-1}(x)}{n!} \right] t^n = \sum_{n=0}^{\infty} \frac{H_{n+1}(x)}{n!} t^n.
$$
Dado que esta igualdad debe cumplirse para todo $t$ en el radio de convergencia, los coeficientes de $t^n$ deben ser idénticos en ambos lados. Igualando los coeficientes y multiplicando toda la ecuación por $n!$, obtenemos:
$$
2x H_n(x) - 2n H_{n-1}(x) = H_{n+1}(x).
$$
Reordenando los términos, llegamos exactamente a la relación de recurrencia solicitada:
$$
2x H_n(x) = 2n H_{n-1}(x) + H_{n+1}(x).
$$
Esta identidad demuestra la conexión entre polinomios de Hermite de órdenes consecutivos, tal como se discute en la Sección 8.5 de Alonso Sepúlveda (p. 338).

\vspace{0.3cm}
\noindent\textbf{Demostración del inciso b):} \\
Derivamos ahora la función generatriz respecto a la variable espacial $x$:
$$
\frac{\partial g}{\partial x} = \frac{\partial}{\partial x} \left( e^{2xt - t^2} \right) = 2t e^{2xt - t^2} = 2t g(x, t).
$$
Derivando la serie de potencias término a término respecto a $x$, obtenemos:
$$
\frac{\partial g}{\partial x} = \sum_{n=0}^{\infty} \frac{H'_n(x)}{n!} t^n.
$$
Igualamos ambas expresiones:
$$
\sum_{n=0}^{\infty} \frac{H'_n(x)}{n!} t^n = 2t \sum_{n=0}^{\infty} \frac{H_n(x)}{n!} t^n = \sum_{n=0}^{\infty} \frac{2 H_n(x)}{n!} t^{n+1}.
$$
Realizamos un cambio de índice en el lado derecho haciendo $m = n + 1$ (es decir, $n = m - 1$), de modo que la suma comienza en $m = 1$:
$$
\sum_{n=0}^{\infty} \frac{H'_n(x)}{n!} t^n = \sum_{m=1}^{\infty} \frac{2 H_{m-1}(x)}{(m-1)!} t^m.
$$
Renombrando el índice mudo $m$ de vuelta a $n$, y usando nuevamente que $\frac{1}{(n-1)!} = \frac{n}{n!}$, podemos escribir el lado derecho comenzando en $n=0$ (pues el término $n=0$ se anula):
$$
\sum_{n=0}^{\infty} \frac{H'_n(x)}{n!} t^n = \sum_{n=0}^{\infty} \frac{2n H_{n-1}(x)}{n!} t^n.
$$
Igualando los coeficientes de $t^n$ y multiplicando por $n!$, obtenemos directamente:
$$
H'_n(x) = 2n H_{n-1}(x).
$$
Esta es una de las identidades fundamentales de los polinomios de Hermite, que relaciona la derivada de un polinomio con el polinomio de orden inmediatamente inferior (Alonso Sepúlveda, Sección 8.5, p. 338).

\vspace{0.3cm}
\noindent\textbf{Demostración del inciso c):} \\
Evaluamos la función generatriz en el punto $x = 0$:
$$
g(0, t) = e^{2(0)t - t^2} = e^{-t^2}.
$$
Por definición de la serie generatriz, esta evaluación también debe ser igual a:
$$
g(0, t) = \sum_{n=0}^{\infty} \frac{H_n(0)}{n!} t^n.
$$
Por otro lado, conocemos la expansión en serie de Maclaurin de la función exponencial $e^u = \sum_{k=0}^{\infty} \frac{u^k}{k!}$. Sustituyendo $u = -t^2$, obtenemos:
$$
e^{-t^2} = \sum_{k=0}^{\infty} \frac{(-t^2)^k}{k!} = \sum_{k=0}^{\infty} \frac{(-1)^k}{k!} t^{2k}.
$$
Igualamos las dos series de potencias en $t$:
$$
\sum_{n=0}^{\infty} \frac{H_n(0)}{n!} t^n = \sum_{k=0}^{\infty} \frac{(-1)^k}{k!} t^{2k}.
$$
Observamos que en el lado derecho solo existen potencias \textbf{pares} de $t$ (es decir, términos de la forma $t^{2k}$). En el lenguaje de series de potencias, esto implica inmediatamente que todos los coeficientes asociados a potencias impares de $t$ en el lado izquierdo deben ser nulos. Por lo tanto, para $n$ impar ($n = 2k + 1$):
$$
\frac{H_{2k+1}(0)}{(2k+1)!} = 0 \quad \Longrightarrow \quad H_{2k+1}(0) = 0.
$$
Esta propiedad refleja la simetría de paridad de los polinomios de Hermite: $H_n(-x) = (-1)^n H_n(x)$. Al evaluar en $x=0$, los polinomios de orden impar (que son funciones impares) deben anularse (Alonso Sepúlveda, Sección 8.5, p. 338).

Para los índices pares, hacemos la sustitución $n = 2k$. Igualando los coeficientes de $t^{2k}$ en ambas series:
$$
\frac{H_{2k}(0)}{(2k)!} = \frac{(-1)^k}{k!}.
$$
Despejando $H_{2k}(0)$, obtenemos la expresión general para los polinomios de Hermite de orden par evaluados en el origen. Renombrando el índice $k$ como $n$ para ajustar la notación al enunciado del problema:
$$
H_{2n}(0) = (-1)^n \frac{(2n)!}{n!}, \quad n = 0, 1, 2, \dots
$$

\vspace{0.3cm}
\noindent\textbf{Demostración del inciso d):} \\
Partimos de la relación obtenida en el inciso b):
$$
H'_n(x) = 2n H_{n-1}(x). \quad (1)
$$
Del inciso a), podemos reescribir la relación de recurrencia despejando $H_{n+1}(x)$:
$$
H_{n+1}(x) = 2x H_n(x) - 2n H_{n-1}(x).
$$
Sustituyendo la ecuación (1) en esta expresión, reemplazamos $2n H_{n-1}(x)$ por $H'_n(x)$:
$$
H_{n+1}(x) = 2x H_n(x) - H'_n(x). \quad (2)
$$
Ahora, derivamos la ecuación (2) con respecto a $x$:
$$
H'_{n+1}(x) = 2 H_n(x) + 2x H'_n(x) - H''_n(x). \quad (3)
$$
Por otro lado, aplicando la relación del inciso b) pero para el índice $n+1$, sabemos que:
$$
H'_{n+1}(x) = 2(n+1) H_n(x) = 2n H_n(x) + 2 H_n(x). \quad (4)
$$
Igualamos las expresiones (3) y (4) para $H'_{n+1}(x)$:
$$
2n H_n(x) + 2 H_n(x) = 2 H_n(x) + 2x H'_n(x) - H''_n(x).
$$
Restamos $2 H_n(x)$ de ambos lados de la ecuación:
$$
2n H_n(x) = 2x H'_n(x) - H''_n(x).
$$
Reordenando todos los términos en un solo lado para igualar a cero, obtenemos la ecuación diferencial de Hermite:
$$
H''_n(x) - 2x H'_n(x) + 2n H_n(x) = 0.
$$
Este desarrollo confirma que los polinomios de Hermite, construidos a partir de la función generatriz y sus relaciones de recurrencia, satisfacen naturalmente esta EDO de segundo orden, la cual es la forma autoadjunta del problema de Sturm-Liouville asociado al oscilador armónico cuántico (Alonso Sepúlveda, Sección 8.5, p. 334 y p. 341).

\vspace{0.3cm}
\noindent\textbf{Demostración del inciso e):} \\
Evaluamos la función generatriz en el valor específico del parámetro $t = 1$. Sustituyendo directamente en la forma cerrada de la función generatriz:
$$
g(x, 1) = e^{2x(1) - (1)^2} = e^{2x - 1} = e^{2x} \cdot e^{-1} = \frac{e^{2x}}{e}.
$$
Por otro lado, evaluando $t = 1$ en la expansión en serie de la función generatriz, obtenemos:
$$
g(x, 1) = \sum_{n=0}^{\infty} \frac{H_n(x)}{n!} (1)^n = \sum_{n=0}^{\infty} \frac{H_n(x)}{n!}.
$$
Igualando ambas expresiones para $g(x, 1)$:
$$
\frac{e^{2x}}{e} = \sum_{n=0}^{\infty} \frac{H_n(x)}{n!}.
$$
Finalmente, multiplicando ambos lados de la ecuación por la constante $e$, obtenemos la identidad solicitada:
$$
e^{2x} = e \sum_{n=0}^{\infty} \frac{H_n(x)}{n!}.
$$
Esta relación es una consecuencia directa y elegante de la definición de la función generatriz, demostrando cómo la evaluación en puntos específicos del parámetro auxiliar $t$ revela identidades funcionales profundas de los polinomios ortogonales (Alonso Sepúlveda, Sección 8.5, p. 338).

\section{Problema 2}
Derivando $n$ veces respecto a $t$ la función generatriz, obtenga la fórmula de Rodrigues:
$$
H_n(x) = (-1)^n e^{x^2} \frac{d^n}{dx^n} \left( e^{-x^2} \right).
$$
Con esta, muestre que los polinomios de Hermite satisfacen la propiedad de paridad:
$$
H_n(-x) = (-1)^n H_n(x).
$$

\subsection{Solución}
Para resolver este problema nos basamos en la teoría de los polinomios de Hermite presentada en el Capítulo 8, Sección 8.5 del texto de Alonso Sepúlveda. La función generatriz de estos polinomios se define como (Alonso Sepúlveda, Sección 8.5, p. 338):
$$
g(x, t) = e^{2xt - t^2} = \sum_{n=0}^{\infty} \frac{H_n(x)}{n!} t^n, \quad \text{para } |t| < \infty.
$$
A partir de esta expresión fundamental, podemos derivar identidades clave mediante diferenciación parcial y evaluación en puntos específicos. A continuación, se detalla el desarrollo analítico paso a paso.

Para acompañar la demostración, se muestran dos gráficos pertinentes. El primero ilustra cómo la fórmula de Rodrigues relaciona los polinomios de Hermite con derivadas sucesivas de la gaussiana $e^{-x^2}$. El segundo compara $H_n(x)$ con $H_n(-x)$ para un caso par y uno impar, haciendo visible la propiedad $H_n(-x)=(-1)^nH_n(x)$.

\begin{center}
\begin{tikzpicture}
\begin{axis}[
    width=0.78\textwidth,
    height=6cm,
    xmin=-3, xmax=3,
    ymin=-1.1, ymax=1.1,
    axis lines=middle,
    grid=both,
    grid style={black!12},
    xlabel={$x$},
    ylabel={valor},
    title={Gaussiana y derivadas asociadas a Rodrigues},
    legend style={font=\scriptsize, cells={anchor=west}, at={(0.02,0.98)}, anchor=north west},
    samples=220,
    domain=-3:3
]
\addplot[blue, thick] {exp(-x^2)};
\addlegendentry{$e^{-x^2}$}
\addplot[red, thick] {-2*x*exp(-x^2)};
\addlegendentry{$\frac{d}{dx}e^{-x^2}$}
\addplot[green!55!black, thick] {(4*x^2-2)*exp(-x^2)};
\addlegendentry{$\frac{d^2}{dx^2}e^{-x^2}$}
\addplot[purple, thick] {(-8*x^3+12*x)*exp(-x^2)};
\addlegendentry{$\frac{d^3}{dx^3}e^{-x^2}$}
\end{axis}
\end{tikzpicture}

\small\textbf{Figura 3.} Primeras derivadas de $e^{-x^2}$. La fórmula de Rodrigues muestra que, al multiplicarlas por $(-1)^n e^{x^2}$, se obtienen los polinomios $H_n(x)$.
\end{center}

\begin{center}
\begin{tikzpicture}
\begin{axis}[
    width=0.78\textwidth,
    height=6cm,
    xmin=-3, xmax=3,
    ymin=-1.15, ymax=1.15,
    axis lines=middle,
    grid=both,
    grid style={black!12},
    xlabel={$x$},
    ylabel={valor normalizado},
    title={Comparación de paridad en $H_2$ y $H_3$},
    legend style={font=\scriptsize, cells={anchor=west}, at={(0.02,0.98)}, anchor=north west},
    samples=220,
    domain=-3:3
]
\addplot[blue, thick] {(4*x^2-2)/34};
\addlegendentry{$H_2(x)/34$}
\addplot[blue, dashed, thick] {(4*x^2-2)/34};
\addlegendentry{$H_2(-x)/34$}
\addplot[red, thick] {(8*x^3-12*x)/180};
\addlegendentry{$H_3(x)/180$}
\addplot[red, dashed, thick] {(-8*x^3+12*x)/180};
\addlegendentry{$H_3(-x)/180$}
\end{axis}
\end{tikzpicture}

\small\textbf{Figura 4.} Visualización de la paridad: $H_2(-x)=H_2(x)$ por ser de orden par, mientras que $H_3(-x)=-H_3(x)$ por ser de orden impar.
\end{center}

\vspace{0.3cm}
\noindent\textbf{Paso 1: Obtención de la fórmula de Rodrigues} \\
Derivamos la función generatriz $n$ veces con respecto al parámetro $t$. Al aplicar el operador $\frac{\partial^n}{\partial t^n}$ a la serie de potencias, obtenemos:
$$
\frac{\partial^n}{\partial t^n} g(x, t) = \sum_{k=n}^{\infty} \frac{H_k(x)}{k!} k(k-1)\cdots(k-n+1) t^{k-n}.
$$
Evaluando esta derivada en $t=0$, todos los términos de la sumatoria se anulan excepto aquel para el cual $k=n$, cuyo coeficiente se simplifica a $n!/n! = 1$. Por lo tanto:
$$
\left. \frac{\partial^n}{\partial t^n} g(x, t) \right|_{t=0} = H_n(x). \quad (1)
$$
Por otro lado, podemos reescribir el exponente de la forma cerrada de la función generatriz completando el cuadrado:
$$
2xt - t^2 = x^2 - (x^2 - 2xt + t^2) = x^2 - (x-t)^2.
$$
Así, la función generatriz se factoriza como:
$$
g(x, t) = e^{x^2} e^{-(x-t)^2}.
$$
Ahora derivamos esta expresión $n$ veces con respecto a $t$. Para facilitar el cálculo, introducimos la variable auxiliar $u = x - t$, de modo que la derivada respecto a $t$ se relaciona con la derivada respecto a $u$ mediante la regla de la cadena: $\frac{\partial}{\partial t} = -\frac{\partial}{\partial u}$. Entonces:
$$
\frac{\partial^n}{\partial t^n} \left[ e^{x^2} e^{-(x-t)^2} \right] = e^{x^2} \left( -\frac{\partial}{\partial u} \right)^n e^{-u^2} = (-1)^n e^{x^2} \frac{\partial^n}{\partial u^n} \left( e^{-u^2} \right).
$$
Evaluamos este resultado en $t=0$, lo que equivale a establecer $u=x$:
$$
\left. \frac{\partial^n}{\partial t^n} g(x, t) \right|_{t=0} = (-1)^n e^{x^2} \left. \frac{d^n}{du^n} \left( e^{-u^2} \right) \right|_{u=x} = (-1)^n e^{x^2} \frac{d^n}{dx^n} \left( e^{-x^2} \right). \quad (2)
$$
Igualando las expresiones (1) y (2), obtenemos directamente la fórmula de Rodrigues para los polinomios de Hermite:
$$
H_n(x) = (-1)^n e^{x^2} \frac{d^n}{dx^n} \left( e^{-x^2} \right).
$$
Esta identidad es fundamental y se encuentra explícitamente listada entre las propiedades algebraicas de los polinomios de Hermite en la Sección 8.5 del texto de Alonso Sepúlveda (p. 338).

\vspace{0.3cm}
\noindent\textbf{Paso 2: Demostración de la propiedad de paridad} \\
Utilizando la fórmula de Rodrigues recién obtenida, evaluamos el polinomio en el punto $-x$:
$$
H_n(-x) = (-1)^n e^{(-x)^2} \frac{d^n}{d(-x)^n} \left( e^{-(-x)^2} \right).
$$
Dado que el cuadrado de un número real es independiente de su signo, $(-x)^2 = x^2$, por lo que los términos exponenciales permanecen invariantes:
$$
H_n(-x) = (-1)^n e^{x^2} \frac{d^n}{d(-x)^n} \left( e^{-x^2} \right).
$$
Ahora, analizamos el operador de derivada. Sabemos que la derivada respecto a $-x$ introduce un signo menos: $\frac{d}{d(-x)} = -\frac{d}{dx}$. Al aplicar este operador $n$ veces consecutivas, obtenemos un factor global de $(-1)^n$:
$$
\frac{d^n}{d(-x)^n} = \left( -\frac{d}{dx} \right)^n = (-1)^n \frac{d^n}{dx^n}.
$$
Sustituyendo esta relación en la expresión para $H_n(-x)$, tenemos:
$$
H_n(-x) = (-1)^n e^{x^2} \left[ (-1)^n \frac{d^n}{dx^n} \left( e^{-x^2} \right) \right].
$$
Agrupando los factores $(-1)^n$, obtenemos $(-1)^{2n}$. Dado que $n$ es un entero no negativo, $(-1)^{2n} = 1$, por lo que:
$$
H_n(-x) = e^{x^2} \frac{d^n}{dx^n} \left( e^{-x^2} \right).
$$
Para recuperar la expresión original de $H_n(x)$, multiplicamos y dividimos por $(-1)^n$ (notando que $1 = (-1)^n \cdot (-1)^n$):
$$
H_n(-x) = (-1)^n \left[ (-1)^n e^{x^2} \frac{d^n}{dx^n} \left( e^{-x^2} \right) \right] = (-1)^n H_n(x).
$$
Esto demuestra rigurosamente que los polinomios de Hermite tienen paridad definida: son funciones pares si $n$ es par, e impares si $n$ es impar. Esta propiedad de simetría es consistente con la discusión general sobre la paridad de las soluciones a ecuaciones diferenciales de Sturm-Liouville con coeficientes pares, como se detalla en la Sección 8.5 de Alonso Sepúlveda, y es crucial para simplificar integrales en problemas de mecánica cuántica, como el cálculo de elementos de matriz en el oscilador armónico.


\section{Problema 3}
Para el oscilador armónico, muestre
$$
\langle \psi_m | x | \psi_n \rangle = \sqrt{\frac{\hbar}{2m\omega}} \left[ \sqrt{n} \delta_{m,n-1} + \sqrt{n+1} \delta_{m,n+1} \right]
$$
¿Qué transiciones dipolares son posibles? (en otras palabras, escriba las reglas de selección dipolares para el oscilador armónico).

\subsection{Solución}
Para resolver este problema nos basamos en el tratamiento algebraico del oscilador armónico cuántico presentado en el Capítulo 8, Sección 8.5.2 del texto de Alonso Sepúlveda. En esta sección se introduce el formalismo de los operadores de creación y aniquilación, los cuales simplifican enormemente el cálculo de elementos de matriz en la base de los estados estacionarios del oscilador.

Antes de realizar el cálculo algebraico, es útil visualizar la idea física: el operador dipolar $\hat{x}$ conecta únicamente estados vecinos. Por eso, en el diagrama de niveles solo aparecen flechas entre niveles adyacentes y, en la matriz de elementos $\langle m|\hat{x}|n\rangle$, los valores no nulos quedan concentrados en las dos diagonales inmediatamente vecinas a la diagonal principal.

\begin{center}
\begin{tikzpicture}[x=1cm,y=0.75cm]
    \foreach \n in {0,1,2,3,4,5} {
        \draw[thick] (0,\n) -- (4.2,\n);
        \node[left] at (0,\n) {$E_{\n}=\hbar\omega(\n+\frac12)$};
        \node[right] at (4.2,\n) {$|\n\rangle$};
    }
    \foreach \n/\m in {0/1,1/2,2/3,3/4,4/5} {
        \draw[-{Latex}, blue!70!black, very thick, bend left=18] (4.75,\n) to node[right, font=\scriptsize] {$\Delta n=+1$} (4.75,\m);
        \draw[-{Latex}, red!70!black, very thick, bend left=18] (-0.55,\m) to node[left, font=\scriptsize] {$\Delta n=-1$} (-0.55,\n);
    }
    \draw[-{Latex}, black!45, dashed, thick] (2.0,0) -- node[pos=0.38, right, font=\scriptsize] {prohibida: $\Delta n=2$} (2.0,2);
    \node[align=center, font=\small] at (2.1,-0.85) {Transiciones dipolares del oscilador armónico};
\end{tikzpicture}

\small\textbf{Figura 5.} Diagrama de niveles de energía. El operador dipolar $\hat{x}\propto\hat{a}+\hat{a}^\dagger$ solo sube o baja un cuanto, por lo que las transiciones permitidas cumplen $\Delta n=\pm1$.
\end{center}

\begin{center}
\begin{tikzpicture}[scale=0.72]
    \foreach \m in {0,...,6} {
        \node[left] at (-0.15,6-\m+0.5) {$m=\m$};
        \node[below, rotate=45, anchor=north east] at (\m+0.5,-0.15) {$n=\m$};
    }
    \foreach \m in {0,...,6} {
        \foreach \n in {0,...,6} {
            \pgfmathtruncatemacro{\diff}{abs(\m-\n)}
            \ifnum\diff=1
                \fill[blue!55] (\n,6-\m) rectangle +(1,1);
                \node[white, font=\scriptsize] at (\n+0.5,6-\m+0.5) {$\neq0$};
            \else
                \fill[black!5] (\n,6-\m) rectangle +(1,1);
            \fi
            \draw[black!35] (\n,6-\m) rectangle +(1,1);
        }
    }
    \node[above] at (3.5,7.25) {Estructura de $\langle m|\hat{x}|n\rangle$};
    \node[below, align=center, font=\small] at (3.5,-1.35) {Solo sobreviven las celdas con $m=n\pm1$};
\end{tikzpicture}

\small\textbf{Figura 6.} Patrón de selección dipolar en forma matricial. Las dos diagonales azules corresponden a los términos $\delta_{m,n-1}$ y $\delta_{m,n+1}$.
\end{center}

\vspace{0.3cm}
\noindent\textbf{Paso 1: Expresión del operador de posición en términos de operadores escalera} \\
El Hamiltoniano del oscilador armónico unidimensional está dado por:
$$
\hat{H} = \frac{\hat{p}^2}{2m} + \frac{1}{2} m \omega^2 \hat{x}^2
$$
Para adimensionalizar el problema, se introduce la variable adimensional $y = x \sqrt{\frac{m\omega}{\hbar}}$. En términos de esta variable, los operadores de aniquilación ($\hat{a}$) y creación ($\hat{a}^\dagger$) se definen como (Alonso Sepúlveda, Sección 8.5.2, p. 344):
$$
\hat{a} = \frac{1}{\sqrt{2}} \left( y + \frac{d}{dy} \right), \quad \hat{a}^\dagger = \frac{1}{\sqrt{2}} \left( y - \frac{d}{dy} \right)
$$
Sumando estas dos expresiones, podemos despejar la variable adimensional $y$:
$$
\hat{a} + \hat{a}^\dagger = \frac{1}{\sqrt{2}} (2y) = \sqrt{2} y \quad \Longrightarrow \quad y = \frac{1}{\sqrt{2}} (\hat{a} + \hat{a}^\dagger)
$$
Dado que $\hat{x} = \sqrt{\frac{\hbar}{m\omega}} y$, el operador de posición $\hat{x}$ puede escribirse en términos de los operadores escalera como:
$$
\hat{x} = \sqrt{\frac{\hbar}{2m\omega}} (\hat{a} + \hat{a}^\dagger)
$$

\vspace{0.3cm}
\noindent\textbf{Paso 2: Acción de los operadores escalera sobre los estados propios} \\
Los estados propios de energía del oscilador armónico, denotados como $|\psi_n\rangle$ o simplemente $|n\rangle$, son autofunciones del operador número $\hat{N} = \hat{a}^\dagger \hat{a}$. La acción de los operadores de aniquilación y creación sobre estos estados está dada por las relaciones fundamentales (Alonso Sepúlveda, Sección 8.5.2, Ecuación 8.72, p. 345):
$$
\hat{a} |n\rangle = \sqrt{n} |n-1\rangle
$$
$$
\hat{a}^\dagger |n\rangle = \sqrt{n+1} |n+1\rangle
$$
Además, recordamos que los estados propios forman una base ortonormal, por lo que satisfacen la condición $\langle m | n \rangle = \delta_{m,n}$.

\vspace{0.3cm}
\noindent\textbf{Paso 3: Cálculo del elemento de matriz $\langle \psi_m | \hat{x} | \psi_n \rangle$} \\
Sustituimos la expresión de $\hat{x}$ en el elemento de matriz solicitado:
$$
\langle \psi_m | \hat{x} | \psi_n \rangle = \sqrt{\frac{\hbar}{2m\omega}} \langle m | (\hat{a} + \hat{a}^\dagger) | n \rangle
$$
Por la linealidad del producto interno, podemos separar esto en dos términos:
$$
\langle \psi_m | \hat{x} | \psi_n \rangle = \sqrt{\frac{\hbar}{2m\omega}} \left( \langle m | \hat{a} | n \rangle + \langle m | \hat{a}^\dagger | n \rangle \right)
$$
Ahora aplicamos la acción de los operadores sobre el ket $|n\rangle$:
$$
\langle m | \hat{a} | n \rangle = \langle m | \left( \sqrt{n} |n-1\rangle \right) = \sqrt{n} \langle m | n-1 \rangle = \sqrt{n} \delta_{m, n-1}
$$
$$
\langle m | \hat{a}^\dagger | n \rangle = \langle m | \left( \sqrt{n+1} |n+1\rangle \right) = \sqrt{n+1} \langle m | n+1 \rangle = \sqrt{n+1} \delta_{m, n+1}
$$
Sustituyendo estos resultados en la expresión del elemento de matriz, obtenemos finalmente:
$$
\langle \psi_m | \hat{x} | \psi_n \rangle = \sqrt{\frac{\hbar}{2m\omega}} \left[ \sqrt{n} \delta_{m, n-1} + \sqrt{n+1} \delta_{m, n+1} \right]
$$
lo cual demuestra la relación solicitada.

\vspace{0.3cm}
\noindent\textbf{Paso 4: Reglas de selección dipolares} \\
En mecánica cuántica, la probabilidad de transición por unidad de tiempo entre un estado inicial $|n\rangle$ y un estado final $|m\rangle$ debido a la interacción con un campo electromagnético (en la aproximación dipolar eléctrica) es proporcional al cuadrado del módulo del elemento de matriz del operador de momento dipolar, $\hat{d} = q\hat{x}$. Es decir, la probabilidad de transición $P_{m \leftarrow n}$ es proporcional a:
$$
P_{m \leftarrow n} \propto |\langle \psi_m | \hat{x} | \psi_n \rangle|^2
$$
Observando la expresión obtenida para el elemento de matriz, vemos que este es diferente de cero únicamente si se cumple alguna de las siguientes condiciones impuestas por las deltas de Kronecker:
\begin{itemize}
    \item $\delta_{m, n-1} \neq 0 \implies m = n - 1$
    \item $\delta_{m, n+1} \neq 0 \implies m = n + 1$
\end{itemize}
Esto significa que el oscilador armónico solo puede realizar transiciones entre niveles de energía adyacentes. En términos del cambio en el número cuántico $\Delta n = m - n$, las únicas transiciones permitidas son aquellas para las cuales:
$$
\Delta n = \pm 1
$$
Esta es la \textbf{regla de selección dipolar} para el oscilador armónico cuántico. Como se discute en la Sección 8.5.2 (p. 345-346) de Alonso Sepúlveda, esto explica por qué el oscilador armónico emite o absorbe radiación de una sola frecuencia fundamental $\omega$, correspondiente a la diferencia de energía $\Delta E = E_{n+1} - E_n = \hbar\omega$, independientemente del nivel de energía en el que se encuentre el sistema. Las transiciones con $\Delta n = \pm 2, \pm 3, \dots$ están prohibidas en la aproximación dipolar eléctrica (aunque pueden ocurrir con probabilidades mucho menores a través de momentos multipolares de orden superior, como el cuadrupolo, donde $\Delta n = \pm 2$).

\section{Problema 4}
Para el oscilador armónico, encuentre las reglas de selección cuadrupolares, es decir, calcule
$$
\langle \psi_m | x^2 | \psi_n \rangle.
$$

\subsection{Solución}
Para resolver este problema nos basamos en el tratamiento algebraico del oscilador armónico cuántico presentado en el Capítulo 8, Sección 8.5.2 del texto de Alonso Sepúlveda. En esta sección se introduce el formalismo de los operadores de creación ($\hat{a}^\dagger$) y aniquilación ($\hat{a}$), los cuales simplifican enormemente el cálculo de elementos de matriz en la base de los estados estacionarios del oscilador, evitando integrales complicadas con polinomios de Hermite.

Antes de iniciar el cálculo, conviene visualizar la estructura física del operador cuadrupolar. Como $\hat{x}^2$ contiene los términos $\hat{a}^2$, $\hat{a}^{\dagger 2}$ y $2\hat{N}+1$, puede conectar un estado consigo mismo o con estados separados por dos cuantos. Además, al ser un operador par, conserva la paridad del estado.

\begin{center}
\begin{tikzpicture}[x=1cm,y=0.72cm]
    \fill[green!8, rounded corners] (-0.35,-0.35) rectangle (2.65,6.55);
    \fill[orange!10, rounded corners] (3.05,-0.35) rectangle (6.05,6.55);
    \node[green!45!black, font=\bfseries\small] at (1.15,6.85) {sector par};
    \node[orange!70!black, font=\bfseries\small] at (4.55,6.85) {sector impar};

    \foreach \n/\y/\sector/\paridad in {0/0/1.15/par,1/1/4.55/impar,2/2/1.15/par,3/3/4.55/impar,4/4/1.15/par,5/5/4.55/impar,6/6/1.15/par} {
        \draw[thick] (\sector-0.85,\y) -- (\sector+0.85,\y);
        \node[left, font=\scriptsize] at (\sector-0.9,\y) {$|\n\rangle$};
        \node[right, font=\scriptsize] at (\sector+0.9,\y) {$E_{\n}$};
    }

    \foreach \a/\b in {0/2,2/4,4/6} {
        \draw[-{Latex}, blue!70!black, very thick, bend left=16] (1.65,\a) to (1.65,\b);
        \draw[-{Latex}, red!70!black, very thick, bend left=16] (0.65,\b) to (0.65,\a);
    }
    \foreach \a/\b in {1/3,3/5} {
        \draw[-{Latex}, blue!70!black, very thick, bend left=16] (5.05,\a) to (5.05,\b);
        \draw[-{Latex}, red!70!black, very thick, bend left=16] (4.05,\b) to (4.05,\a);
    }

    \foreach \n/\sector in {0/1.15,1/4.55,2/1.15,3/4.55,4/1.15,5/4.55,6/1.15} {
        \draw[black!55, thick, loop right, min distance=8mm] (\sector+0.35,\n) to (\sector+0.35,\n);
    }

    \node[draw, rounded corners, fill=white, align=left, font=\scriptsize] at (3.1,-1.0) {
        $\hat{a}^{\dagger 2}:\ \Delta n=+2$\quad $\hat{a}^{2}:\ \Delta n=-2$\\
        $2\hat{N}+1:\ \Delta n=0$\quad $\hat{x}^2$ conserva paridad
    };
\end{tikzpicture}

\small\textbf{Figura 7.} Diagrama de niveles para las transiciones cuadrupolares. Las flechas azules y rojas conectan niveles separados por dos cuantos, mientras que los lazos representan los términos diagonales $m=n$.
\end{center}

\begin{center}
\begin{tikzpicture}[scale=0.68]
    \foreach \m in {0,...,7} {
        \node[left, font=\scriptsize] at (-0.15,7-\m+0.5) {$m=\m$};
        \node[below, rotate=45, anchor=north east, font=\scriptsize] at (\m+0.5,-0.15) {$n=\m$};
    }
    \foreach \m in {0,...,7} {
        \foreach \n in {0,...,7} {
            \pgfmathtruncatemacro{\diff}{\m-\n}
            \pgfmathtruncatemacro{\adiff}{abs(\m-\n)}
            \ifnum\adiff=0
                \fill[purple!60] (\n,7-\m) rectangle +(1,1);
                \node[white, font=\tiny] at (\n+0.5,7-\m+0.5) {$2n+1$};
            \else
                \ifnum\diff=2
                    \fill[blue!58] (\n,7-\m) rectangle +(1,1);
                    \node[white, font=\tiny] at (\n+0.5,7-\m+0.5) {$+2$};
                \else
                    \ifnum\diff=-2
                        \fill[red!58] (\n,7-\m) rectangle +(1,1);
                        \node[white, font=\tiny] at (\n+0.5,7-\m+0.5) {$-2$};
                    \else
                        \fill[black!4] (\n,7-\m) rectangle +(1,1);
                    \fi
                \fi
            \fi
            \draw[black!30] (\n,7-\m) rectangle +(1,1);
        }
    }
    \draw[decorate, decoration={brace, amplitude=5pt}, thick] (8.15,8) -- node[right=6pt, align=left, font=\scriptsize] {solo tres\\diagonales} (8.15,0);
    \node[above, align=center] at (4,8.25) {Patrón de $\langle m|\hat{x}^2|n\rangle$};
    \node[draw, rounded corners, fill=white, align=left, font=\scriptsize] at (4,-1.85) {
        \textcolor{purple!70!black}{diagonal: $m=n$}\quad
        \textcolor{blue!70!black}{$m=n+2$}\quad
        \textcolor{red!70!black}{$m=n-2$}
    };
\end{tikzpicture}

\small\textbf{Figura 8.} Mapa matricial de los elementos cuadrupolares no nulos. La diagonal central proviene de $2\hat{N}+1$, y las diagonales secundarias de $\hat{a}^{\dagger 2}$ y $\hat{a}^{2}$.
\end{center}

\vspace{0.3cm}
\noindent\textbf{Paso 1: Expresión del operador de posición al cuadrado en términos de operadores escalera} \\
El operador de posición $\hat{x}$ se puede escribir en términos de los operadores de aniquilación y creación como (Alonso Sepúlveda, Sección 8.5.2, p. 344):
$$
\hat{x} = \sqrt{\frac{\hbar}{2m\omega}} (\hat{a} + \hat{a}^\dagger)
$$
Para encontrar el elemento de matriz cuadrupolar, necesitamos elevar al cuadrado este operador:
$$
\hat{x}^2 = \frac{\hbar}{2m\omega} (\hat{a} + \hat{a}^\dagger)^2 = \frac{\hbar}{2m\omega} (\hat{a}^2 + \hat{a}^{\dagger 2} + \hat{a}\hat{a}^\dagger + \hat{a}^\dagger\hat{a})
$$
Utilizamos la relación de conmutación fundamental de los operadores escalera, $[\hat{a}, \hat{a}^\dagger] = 1$, la cual implica que $\hat{a}\hat{a}^\dagger = \hat{a}^\dagger\hat{a} + 1$. Sustituyendo esto en la expresión anterior, obtenemos una forma ordenada (con el operador número $\hat{N} = \hat{a}^\dagger\hat{a}$):
$$
\hat{x}^2 = \frac{\hbar}{2m\omega} (\hat{a}^2 + \hat{a}^{\dagger 2} + 2\hat{a}^\dagger\hat{a} + 1)
$$

\vspace{0.3cm}
\noindent\textbf{Paso 2: Acción de los operadores sobre los estados propios} \\
Los estados propios de energía del oscilador armónico, denotados como $|\psi_n\rangle$ o simplemente $|n\rangle$, satisfacen las siguientes relaciones de acción de los operadores escalera (Alonso Sepúlveda, Sección 8.5.2, Ecuación 8.72, p. 344):
$$
\hat{a} |n\rangle = \sqrt{n} |n-1\rangle
$$
$$
\hat{a}^\dagger |n\rangle = \sqrt{n+1} |n+1\rangle
$$
Además, recordamos que los estados propios forman una base ortonormal, por lo que $\langle m | n \rangle = \delta_{m,n}$.

\vspace{0.3cm}
\noindent\textbf{Paso 3: Cálculo del elemento de matriz $\langle \psi_m | \hat{x}^2 | \psi_n \rangle$} \\
Sustituimos la expresión de $\hat{x}^2$ en el elemento de matriz solicitado y aplicamos la linealidad del producto interno:
$$
\langle m | \hat{x}^2 | n \rangle = \frac{\hbar}{2m\omega} \langle m | \left( \hat{a}^2 + \hat{a}^{\dagger 2} + 2\hat{a}^\dagger\hat{a} + 1 \right) | n \rangle
$$
Evaluamos cada término por separado:

1. \textbf{Término $\hat{a}^2$}:
$$
\langle m | \hat{a}^2 | n \rangle = \langle m | \hat{a} (\sqrt{n} |n-1\rangle) = \sqrt{n} \langle m | \sqrt{n-1} |n-2\rangle = \sqrt{n(n-1)} \langle m | n-2 \rangle = \sqrt{n(n-1)} \delta_{m, n-2}
$$

2. \textbf{Término $\hat{a}^{\dagger 2}$}:
\begin{align*}
\langle m | \hat{a}^{\dagger 2} | n \rangle
&= \langle m | \hat{a}^\dagger (\sqrt{n+1} |n+1\rangle) \\
&= \sqrt{n+1}\,\langle m | \sqrt{n+2} |n+2\rangle \\
&= \sqrt{(n+1)(n+2)}\,\langle m | n+2 \rangle \\
&= \sqrt{(n+1)(n+2)}\,\delta_{m, n+2}.
\end{align*}

3. \textbf{Término $2\hat{a}^\dagger\hat{a}$}:
$$
\langle m | 2\hat{a}^\dagger\hat{a} | n \rangle = 2n \langle m | n \rangle = 2n \delta_{m, n}
$$

4. \textbf{Término $1$ (Identidad)}:
$$
\langle m | 1 | n \rangle = \langle m | n \rangle = \delta_{m, n}
$$

Sumando todos los términos, obtenemos la expresión completa para el elemento de matriz cuadrupolar:
$$
\langle \psi_m | \hat{x}^2 | \psi_n \rangle = \frac{\hbar}{2m\omega} \left[ \sqrt{n(n-1)} \delta_{m, n-2} + \sqrt{(n+1)(n+2)} \delta_{m, n+2} + (2n+1) \delta_{m, n} \right]
$$

\vspace{0.3cm}
\noindent\textbf{Paso 4: Reglas de selección cuadrupolares} \\
En mecánica cuántica, la probabilidad de una transición entre un estado inicial $|n\rangle$ y un estado final $|m\rangle$ mediada por un operador (en este caso, el momento cuadrupolar, proporcional a $\hat{x}^2$) es proporcional al cuadrado del módulo del elemento de matriz correspondiente. 

Observando la expresión obtenida, el elemento de matriz es diferente de cero únicamente si se cumple alguna de las siguientes condiciones impuestas por las deltas de Kronecker:
\begin{itemize}
    \item $\delta_{m, n-2} \neq 0 \implies m = n - 2$
    \item $\delta_{m, n+2} \neq 0 \implies m = n + 2$
    \item $\delta_{m, n} \neq 0 \implies m = n$
\end{itemize}

Esto significa que el oscilador armónico solo puede realizar transiciones cuadrupolares entre niveles de energía que difieran en $0$ o $\pm 2$ cuantos de energía. En términos del cambio en el número cuántico $\Delta n = m - n$, las reglas de selección cuadrupolares son:
$$
\Delta n = 0, \pm 2
$$

\vspace{0.3cm}
\noindent\textbf{Discusión física y paridad} \\
Este resultado es completamente consistente con las propiedades de simetría del sistema. El operador $\hat{x}^2$ es una función par bajo inversión espacial ($x \to -x$). Los estados propios del oscilador armónico tienen paridad definida: $P|n\rangle = (-1)^n |n\rangle$. Para que un elemento de matriz de un operador par sea no nulo, los estados inicial y final deben tener la misma paridad, es decir, $(-1)^m = (-1)^n$, lo cual exige que la diferencia $m - n$ sea un número par. 

Mientras que las transiciones dipolares (gobernadas por el operador $\hat{x}$, que es impar) requieren $\Delta n = \pm 1$ (cambio de paridad), las transiciones cuadrupolares (gobernadas por $\hat{x}^2$, que es par) requieren $\Delta n = 0, \pm 2$ (conservación de paridad), tal como se discute en el contexto de las reglas de selección en la Sección 8.5.2 (p. 345-346) de Alonso Sepúlveda. Estas transiciones son de orden superior y, por lo tanto, tienen probabilidades de ocurrencia mucho menores que las dipolares, pero son fundamentales en espectroscopía cuando las transiciones dipolares están prohibidas por las reglas de selección.


\section{Problema 5}
Demuestre la relación de ortogonalidad de los polinomios de Hermite:
$$
\int_{-\infty}^{\infty} e^{-x^2} H_m(x) H_n(x) \, dx = \sqrt{\pi} \, 2^n \, n! \, \delta_{mn}.
$$

\subsection{Solución}
Para abordar este problema nos basamos en el desarrollo formal de la ortogonalidad de los polinomios de Hermite presentado en el Capítulo 8, Sección 8.5 (pp. 339--340) del texto de Alonso Sepúlveda. Esta propiedad es fundamental, ya que garantiza que los polinomios de Hermite forman una base ortogonal en el espacio de Hilbert $\mathcal{L}^2(\mathbb{R})$ con peso $w(x) = e^{-x^2}$, lo cual es la base matemática para la cuantización del oscilador armónico y para la expansión de funciones en series de Hermite. A continuación, se detalla la demostración analítica paso a paso.

La idea central de la ortogonalidad puede verse geométricamente si se consideran las funciones de Hermite ponderadas por la gaussiana: cada modo tiene un patrón de nodos distinto y, al multiplicar dos modos diferentes e integrar en todo el eje real, las áreas positivas y negativas se cancelan exactamente. En cambio, el producto de un modo consigo mismo siempre produce una contribución positiva que da la norma $\sqrt{\pi}\,2^n n!$.

\begin{center}
\begin{tikzpicture}
\begin{axis}[
    width=0.58\textwidth,
    height=6.2cm,
    xmin=-3.2, xmax=3.2,
    ymin=-1.25, ymax=1.25,
    axis lines=middle,
    grid=both,
    grid style={black!10},
    xlabel={$x$},
    ylabel={modo ponderado},
    title={Modos gaussianos de Hermite},
    legend style={font=\scriptsize, cells={anchor=west}, at={(0.02,0.98)}, anchor=north west},
    samples=240,
    domain=-3.2:3.2
]
\addplot[blue!75!black, thick] {exp(-x^2/2)};
\addlegendentry{$\Phi_0$}
\addplot[red!75!black, thick] {sqrt(2)*x*exp(-x^2/2)};
\addlegendentry{$\Phi_1$}
\addplot[green!55!black, thick] {((2*x^2-1)/sqrt(2))*exp(-x^2/2)};
\addlegendentry{$\Phi_2$}
\addplot[purple!75!black, thick] {((2*x^3-3*x)/sqrt(3))*exp(-x^2/2)};
\addlegendentry{$\Phi_3$}
\addplot[black!45, dashed, thick] {0};
\end{axis}


\end{tikzpicture}

\small\textbf{Figura 9.} Visualización de la ortogonalidad de Hermite. Se muestran modos gaussianos con paridades y nodos distintos; al multiplicar modos diferentes e integrar sobre todo el eje real, las contribuciones positivas y negativas se cancelan.
\end{center}

\vspace{0.3cm}
\noindent\textbf{Paso 1: Planteamiento de la integral de ortogonalidad} \\
Consideremos la integral general que define el producto interno con peso $e^{-x^2}$ entre dos polinomios de Hermite de órdenes $m$ y $n$:
$$
I_{mn} = \int_{-\infty}^{\infty} e^{-x^2} H_m(x) H_n(x) \, dx.
$$
Para evaluar esta expresión, sustituiremos uno de los polinomios mediante la fórmula de Rodrigues, la cual establecimos en el problema anterior (Alonso Sepúlveda, Sección 8.5, p. 338):
$$
H_n(x) = (-1)^n e^{x^2} \frac{d^n}{dx^n} \left( e^{-x^2} \right).
$$
Insertando esta expresión en $I_{mn}$, los factores exponenciales $e^{x^2}$ y $e^{-x^2}$ se cancelan algebraicamente, simplificando notablemente el integrando:
$$
I_{mn} = (-1)^n \int_{-\infty}^{\infty} H_m(x) \frac{d^n}{dx^n} \left( e^{-x^2} \right) \, dx. \quad (1)
$$

\vspace{0.3cm}
\noindent\textbf{Paso 2: Integración por partes $n$ veces} \\
Aplicaremos el método de integración por partes de forma iterada $n$ veces sobre la ecuación (1). Recordemos la fórmula estándar:
$$
\int u \, dv = uv - \int v \, du.
$$
Tomamos $u = H_m(x)$ y $dv = \frac{d^n}{dx^n}(e^{-x^2}) dx$. En cada paso, derivaremos el polinomio $H_m(x)$ e integraremos la derivada de orden superior de la gaussiana. Tras la primera integración por partes:
$$
I_{mn} = (-1)^n \left[ H_m(x) \frac{d^{n-1}}{dx^{n-1}}(e^{-x^2}) \right]_{-\infty}^{\infty} - (-1)^n \int_{-\infty}^{\infty} H'_m(x) \frac{d^{n-1}}{dx^{n-1}}(e^{-x^2}) \, dx.
$$
El término evaluado en los límites $\pm \infty$ se anula estrictamente. Esto se debe a que la función $e^{-x^2}$ y todas sus derivadas decaen exponencialmente a cero cuando $|x| \to \infty$, mientras que $H_m(x)$ crece solo como un polinomio de grado $m$. Matemáticamente, para cualquier entero $k \geq 0$:
$$
\lim_{x \to \pm \infty} x^k e^{-x^2} = 0 \quad \text{y} \quad \lim_{x \to \pm \infty} \frac{d^k}{dx^k}(e^{-x^2}) = 0.
$$
Por lo tanto, todos los términos de frontera desaparecen en cada aplicación de integración por partes. Repitiendo el proceso $n$ veces consecutivas, obtenemos:
$$
I_{mn} = (-1)^n (-1)^n \int_{-\infty}^{\infty} \frac{d^n}{dx^n} \left[ H_m(x) \right] e^{-x^2} \, dx = \int_{-\infty}^{\infty} e^{-x^2} \frac{d^n H_m(x)}{dx^n} \, dx. \quad (2)
$$
El factor $(-1)^{2n} = 1$ se ha simplificado.

\vspace{0.3cm}
\noindent\textbf{Paso 3: Análisis por casos según los grados $m$ y $n$} \\
La expresión (2) depende crucialmente de la relación entre los órdenes $m$ y $n$, dado que $H_m(x)$ es un polinomio de grado exactamente $m$. Analizamos tres situaciones:

\begin{itemize}
    \item \textbf{Caso 1: $m < n$}. Si el grado del polinomio $H_m(x)$ es menor que el número de derivadas $n$, entonces la derivada $n$-ésima es idénticamente cero:
    $$
    \frac{d^n}{dx^n} H_m(x) = 0 \quad \text{para } m < n.
    $$
    Sustituyendo en (2), obtenemos inmediatamente $I_{mn} = 0$.

    \item \textbf{Caso 2: $m > n$}. Por la simetría del producto interno, $I_{mn} = I_{nm}$. Si intercambiamos los roles de $m$ y $n$ en la demostración anterior, el mismo argumento de integración por partes $m$ veces lleva a una integral que contiene $\frac{d^m}{dx^m} H_n(x)$, la cual es cero porque $n < m$. Por lo tanto, $I_{mn} = 0$ también en este caso.

    \item \textbf{Caso 3: $m = n$}. Este es el caso donde la integral no se anula. La derivada $n$-ésima de un polinomio de grado $n$ es una constante igual al coeficiente del término de mayor grado multiplicado por $n!$. Según las propiedades de los polinomios de Hermite (Alonso Sepúlveda, Sección 8.5, p. 338), el coeficiente del término de mayor potencia $x^n$ en $H_n(x)$ es $2^n$. Es decir:
    $$
    H_n(x) = 2^n x^n + \text{términos de orden inferior}.
    $$
    Al derivar $n$ veces, todos los términos de orden inferior desaparecen y queda:
    $$
    \frac{d^n}{dx^n} H_n(x) = 2^n \frac{d^n}{dx^n} (x^n) = 2^n \cdot n!.
    $$
\end{itemize}

\vspace{0.3cm}
\noindent\textbf{Paso 4: Evaluación final y normalización} \\
Sustituyendo el resultado del caso $m=n$ en la ecuación (2), la constante $2^n n!$ puede sacarse fuera de la integral:
$$
I_{nn} = 2^n n! \int_{-\infty}^{\infty} e^{-x^2} \, dx.
$$
La integral restante es la clásica integral gaussiana, cuyo valor es bien conocido y se encuentra tabulado en la literatura matemática estándar (Alonso Sepúlveda, Apéndices y Sección 8.5, p. 340):
$$
\int_{-\infty}^{\infty} e^{-x^2} \, dx = \sqrt{\pi}.
$$
Por lo tanto:
$$
I_{nn} = \sqrt{\pi} \, 2^n \, n!.
$$

Combinando los tres casos analizados mediante el uso de la delta de Kronecker $\delta_{mn}$ (que vale 1 si $m=n$ y 0 en caso contrario), obtenemos la relación de ortogonalidad completa:
$$
\int_{-\infty}^{\infty} e^{-x^2} H_m(x) H_n(x) \, dx = \sqrt{\pi} \, 2^n \, n! \, \delta_{mn}.
$$

\vspace{0.3cm}
\noindent\textbf{Discusión y relevancia física} \\
Este resultado demuestra que los polinomios de Hermite son ortogonales bajo el peso $e^{-x^2}$ en el dominio $(-\infty, \infty)$. La presencia del factor de normalización $\sqrt{\pi} 2^n n!$ indica que los polinomios no están ortonormalizados, sino simplemente ortogonales. En el contexto de la mecánica cuántica, esta relación se utiliza directamente para normalizar las autofunciones del oscilador armónico unidimensional:
$$
\psi_n(x) = \left( \frac{m\omega}{\pi \hbar} \right)^{1/4} \frac{1}{\sqrt{2^n n!}} H_n\left( \sqrt{\frac{m\omega}{\hbar}} x \right) e^{-\frac{m\omega}{2\hbar} x^2},
$$
garantizando que $\int_{-\infty}^{\infty} |\psi_n(x)|^2 dx = 1$. La demostración presentada sigue rigurosamente la línea de razonamiento expuesta en la Sección 8.5 del texto de Alonso Sepúlveda, donde se enfatiza que la ortogonalidad es consecuencia directa de la estructura de Sturm-Liouville de la ecuación diferencial de Hermite y de la fórmula de Rodrigues.


\section{Problema 6}
Muestre las siguientes relaciones:
$$
\sin(2x) = \frac{1}{e} \sum_{n=0}^{\infty} \frac{(-1)^n}{(2n+1)!} H_{2n+1}(x), \quad \cos(2x) = \frac{1}{e} \sum_{n=0}^{\infty} \frac{(-1)^n}{(2n)!} H_{2n}(x).
$$

\subsection{Solución}
Para resolver este problema nos basamos en la teoría de los polinomios de Hermite presentada en el Capítulo 8, Sección 8.5 del texto de Alonso Sepúlveda. La herramienta fundamental para esta demostración es la función generatriz de los polinomios de Hermite, la cual se define como (Alonso Sepúlveda, Sección 8.5, p. 338):
$$
g(x, t) = e^{2xt - t^2} = \sum_{n=0}^{\infty} \frac{H_n(x)}{n!} t^n, \quad \text{para todo } t \in \mathbb{C}.
$$
Esta serie de potencias converge absolutamente para cualquier valor complejo de $t$, lo que nos permite evaluar la función generatriz en puntos del plano complejo, como $t = i$ (donde $i$ es la unidad imaginaria, $i^2 = -1$), sin perder validez matemática. A continuación, se detalla el desarrollo analítico paso a paso.

Antes de realizar la demostración algebraica, es útil comparar gráficamente las funciones exactas $\sin(2x)$ y $\cos(2x)$ con las sumas truncadas de Hermite. En la figura siguiente se usan los términos hasta $H_9(x)$ para el seno y hasta $H_8(x)$ para el coseno, lo cual permite observar cómo la separación entre términos impares y pares reproduce, respectivamente, la simetría impar del seno y la simetría par del coseno.

\begin{center}
\begin{tikzpicture}
\begin{groupplot}[
    group style={group size=2 by 1, horizontal sep=1.2cm},
    width=0.47\textwidth,
    height=5.8cm,
    xmin=-2.2, xmax=2.2,
    ymin=-1.25, ymax=1.25,
    axis lines=middle,
    grid=both,
    grid style={black!12},
    xlabel={$x$},
    samples=260,
    domain=-2.2:2.2,
    legend style={font=\scriptsize, cells={anchor=west}, at={(0.03,0.97)}, anchor=north west}
]
\nextgroupplot[
    title={Serie impar para $\sin(2x)$},
    ylabel={valor}
]
\addplot[blue, thick] {sin(deg(2*x))};
\addlegendentry{$\sin(2x)$}
\addplot[red, dashed, thick] {exp(-1)*(2*x - (8*x^3-12*x)/6 + (32*x^5-160*x^3+120*x)/120 - (128*x^7-1344*x^5+3360*x^3-1680*x)/5040 + (512*x^9-9216*x^7+48384*x^5-80640*x^3+30240*x)/362880)};
\addlegendentry{suma hasta $H_9$}

\nextgroupplot[
    title={Serie par para $\cos(2x)$}
]
\addplot[blue, thick] {cos(deg(2*x))};
\addlegendentry{$\cos(2x)$}
\addplot[red, dashed, thick] {exp(-1)*(1 - (4*x^2-2)/2 + (16*x^4-48*x^2+12)/24 - (64*x^6-480*x^4+720*x^2-120)/720 + (256*x^8-3584*x^6+13440*x^4-13440*x^2+1680)/40320)};
\addlegendentry{suma hasta $H_8$}
\end{groupplot}
\end{tikzpicture}

\small\textbf{Figura 10.} Comparación entre las funciones trigonométricas exactas y las sumas parciales de Hermite obtenidas al evaluar la función generatriz en $t=i$. Los términos impares aproximan al seno y los términos pares aproximan al coseno.
\end{center}

\vspace{0.3cm}
\noindent\textbf{Paso 1: Evaluación de la forma cerrada de la función generatriz en $t = i$} \\
Sustituimos $t = i$ directamente en la expresión exponencial de la función generatriz:
$$
g(x, i) = e^{2x(i) - (i)^2}.
$$
Dado que $i^2 = -1$, el exponente se simplifica de la siguiente manera:
$$
2x(i) - (i)^2 = 2ix - (-1) = 1 + 2ix.
$$
Por lo tanto, la función generatriz evaluada en $t = i$ se convierte en:
$$
g(x, i) = e^{1 + 2ix} = e^1 \cdot e^{2ix} = e \cdot e^{2ix}.
$$
Aplicando la fórmula de Euler, $e^{i\theta} = \cos(\theta) + i\sin(\theta)$, con $\theta = 2x$, podemos separar esta expresión en sus partes real e imaginaria:
$$
g(x, i) = e \left( \cos(2x) + i \sin(2x) \right). \quad (1)
$$

\vspace{0.3cm}
\noindent\textbf{Paso 2: Evaluación de la serie de potencias en $t = i$} \\
Ahora sustituimos $t = i$ en el lado derecho de la definición de la función generatriz (la serie infinita):
$$
g(x, i) = \sum_{n=0}^{\infty} \frac{H_n(x)}{n!} i^n.
$$
Para analizar esta suma, es conveniente separarla en dos series: una que contenga los términos con índices pares ($n = 2k$) y otra con los términos de índices impares ($n = 2k + 1$), donde $k = 0, 1, 2, \dots$. 

Analicemos las potencias de $i$:
\begin{itemize}
    \item Para $n = 2k$ (par): $i^{2k} = (i^2)^k = (-1)^k$.
    \item Para $n = 2k + 1$ (impar): $i^{2k+1} = i \cdot i^{2k} = i(-1)^k$.
\end{itemize}

Sustituyendo estas identidades en la serie, obtenemos:
$$
g(x, i) = \sum_{k=0}^{\infty} \frac{H_{2k}(x)}{(2k)!} (-1)^k + \sum_{k=0}^{\infty} \frac{H_{2k+1}(x)}{(2k+1)!} i(-1)^k.
$$
Factorizando la unidad imaginaria $i$ en la segunda sumatoria, podemos agrupar las partes real e imaginaria de la serie:
$$
g(x, i) = \left[ \sum_{k=0}^{\infty} \frac{(-1)^k}{(2k)!} H_{2k}(x) \right] + i \left[ \sum_{k=0}^{\infty} \frac{(-1)^k}{(2k+1)!} H_{2k+1}(x) \right]. \quad (2)
$$

\vspace{0.3cm}
\noindent\textbf{Paso 3: Igualación de partes reales e imaginarias} \\
Dado que las expresiones (1) y (2) representan la misma función $g(x, i)$, sus partes reales deben ser iguales entre sí, y sus partes imaginarias también deben ser iguales entre sí. 

\textbf{Igualando las partes reales:}
$$
e \cos(2x) = \sum_{k=0}^{\infty} \frac{(-1)^k}{(2k)!} H_{2k}(x).
$$
Dividiendo ambos lados por la constante $e$, obtenemos la primera relación solicitada:
$$
\cos(2x) = \frac{1}{e} \sum_{k=0}^{\infty} \frac{(-1)^k}{(2k)!} H_{2k}(x).
$$

\textbf{Igualando las partes imaginarias:}
$$
e \sin(2x) = \sum_{k=0}^{\infty} \frac{(-1)^k}{(2k+1)!} H_{2k+1}(x).
$$
Dividiendo nuevamente ambos lados por $e$, obtenemos la segunda relación solicitada:
$$
\sin(2x) = \frac{1}{e} \sum_{k=0}^{\infty} \frac{(-1)^k}{(2k+1)!} H_{2k+1}(x).
$$

\vspace{0.3cm}
\noindent\textbf{Discusión y relevancia} \\
Este elegante resultado demuestra cómo la función generatriz, al ser evaluada en el plano complejo, actúa como un puente directo entre las funciones trigonométricas y los polinomios de Hermite. La separación en partes par e impar de la serie refleja la propiedad de paridad de los polinomios de Hermite: $H_n(-x) = (-1)^n H_n(x)$ (Alonso Sepúlveda, Sección 8.5, p. 338). Los polinomios de orden par $H_{2k}(x)$ son funciones pares y se asocian naturalmente con el coseno (que también es par), mientras que los polinomios de orden impar $H_{2k+1}(x)$ son funciones impares y se asocian con el seno (que es impar). 

Finalmente, renombrando el índice mudo de sumatoria $k$ como $n$ para coincidir exactamente con la notación del enunciado del problema, concluimos la demostración de manera rigurosa:
$$
\sin(2x) = \frac{1}{e} \sum_{n=0}^{\infty} \frac{(-1)^n}{(2n+1)!} H_{2n+1}(x), \quad \cos(2x) = \frac{1}{e} \sum_{n=0}^{\infty} \frac{(-1)^n}{(2n)!} H_{2n}(x).
$$
Esto completa la demostración analítica solicitada, fundamentada en las propiedades de la función generatriz expuestas en el texto de referencia.

\section{Problema 7}
Expandir las funciones
$$
f(x) = x^{2r}, \quad g(x) = x^{2r+1}
$$
en términos de los polinomios de Hermite.

\subsection{Solución}
Para resolver este problema nos basamos en la teoría de los polinomios de Hermite presentada en el Capítulo 8, Sección 8.5 del texto de Alonso Sepúlveda. Cualquier función $f(x)$ de cuadrado integrable con peso $e^{-x^2}$ en el intervalo $(-\infty, \infty)$ puede expandirse en una serie de polinomios de Hermite de la forma (Alonso Sepúlveda, Sección 8.5, p. 339):
$$
f(x) = \sum_{n=0}^{\infty} c_n H_n(x)
$$
donde los coeficientes $c_n$ se determinan utilizando la relación de ortogonalidad de los polinomios de Hermite:
$$
\int_{-\infty}^{\infty} e^{-x^2} H_m(x) H_n(x) \, dx = \sqrt{\pi} \, 2^n \, n! \, \delta_{mn}
$$
Multiplicando la expansión por $e^{-x^2} H_m(x)$ e integrando, obtenemos la fórmula explícita para los coeficientes:
$$
c_n = \frac{1}{\sqrt{\pi} \, 2^n \, n!} \int_{-\infty}^{\infty} e^{-x^2} f(x) H_n(x) \, dx
$$
A continuación, aplicamos este método analítico a las funciones $f(x) = x^{2r}$ y $g(x) = x^{2r+1}$.


Antes de desarrollar los coeficientes en forma general, conviene visualizar la idea de la expansión finita. La figura siguiente muestra dos ejemplos particulares: $x^4$ se reconstruye con polinomios pares $H_0$, $H_2$ y $H_4$, mientras que $x^5$ se reconstruye con polinomios impares $H_1$, $H_3$ y $H_5$. Esto ilustra que la paridad del monomio determina qué familia de polinomios de Hermite aparece en la suma.

\begin{center}
\begin{tikzpicture}
\begin{groupplot}[
    group style={group size=2 by 1, horizontal sep=1.2cm},
    width=0.47\textwidth,
    height=5.8cm,
    xmin=-2, xmax=2,
    axis lines=middle,
    grid=both,
    grid style={black!12},
    xlabel={$x$},
    samples=240,
    domain=-2:2,
    legend style={font=\scriptsize, cells={anchor=west}, at={(0.03,0.97)}, anchor=north west}
]
\nextgroupplot[
    title={Expansión par de $x^4$},
    ymin=-0.5, ymax=16.5,
    ylabel={valor}
]
\addplot[blue, thick] {x^4};
\addlegendentry{$x^4$}
\addplot[red, dashed, thick] {3/4 + (3/4)*(4*x^2-2) + (1/16)*(16*x^4-48*x^2+12)};
\addlegendentry{$\frac{3}{4}H_0+\frac{3}{4}H_2+\frac{1}{16}H_4$}

\nextgroupplot[
    title={Expansión impar de $x^5$},
    ymin=-33, ymax=33
]
\addplot[blue, thick] {x^5};
\addlegendentry{$x^5$}
\addplot[red, dashed, thick] {(15/8)*(2*x) + (5/16)*(8*x^3-12*x) + (1/32)*(32*x^5-160*x^3+120*x)};
\addlegendentry{$\frac{15}{8}H_1+\frac{5}{16}H_3+\frac{1}{32}H_5$}
\end{groupplot}
\end{tikzpicture}

\small\textbf{Figura 11.} Ejemplos de expansiones finitas de monomios en polinomios de Hermite. Los monomios pares usan únicamente $H_{2n}(x)$ y los monomios impares únicamente $H_{2n+1}(x)$.
\end{center}

\vspace{0.3cm}
\noindent\textbf{Caso 1: Expansión de $f(x) = x^{2r}$} \\
Los coeficientes de la expansión están dados por:
$$
c_n = \frac{1}{\sqrt{\pi} \, 2^n \, n!} \int_{-\infty}^{\infty} e^{-x^2} x^{2r} H_n(x) \, dx
$$
Dado que $x^{2r}$ es una función par y los polinomios de Hermite tienen paridad definida $H_n(-x) = (-1)^n H_n(x)$, el integrando es una función impar cuando $n$ es impar. Por lo tanto, la integral se anula para $n$ impar, es decir, $c_{2k+1} = 0$. 

Para $n$ par, hacemos $n = 2m$ (con $m$ entero no negativo). Utilizamos la fórmula de Rodrigues para los polinomios de Hermite (Alonso Sepúlveda, Sección 8.5, p. 338):
$$
H_{2m}(x) = (-1)^{2m} e^{x^2} \frac{d^{2m}}{dx^{2m}} \left( e^{-x^2} \right) = e^{x^2} \frac{d^{2m}}{dx^{2m}} \left( e^{-x^2} \right)
$$
Sustituyendo esto en la integral para $c_{2m}$:
$$
c_{2m} = \frac{1}{\sqrt{\pi} \, 2^{2m} \, (2m)!} \int_{-\infty}^{\infty} x^{2r} \frac{d^{2m}}{dx^{2m}} \left( e^{-x^2} \right) \, dx
$$
Aplicamos integración por partes $2m$ veces. Dado que $e^{-x^2}$ y todas sus derivadas tienden a cero cuando $x \to \pm \infty$, los términos de frontera se anulan. Cada integración por partes introduce un factor de $(-1)$, por lo que después de $2m$ pasos obtenemos:
$$
c_{2m} = \frac{(-1)^{2m}}{\sqrt{\pi} \, 2^{2m} \, (2m)!} \int_{-\infty}^{\infty} e^{-x^2} \frac{d^{2m}}{dx^{2m}} \left( x^{2r} \right) \, dx
$$
La derivada $2m$-ésima de $x^{2r}$ es:
$$
\frac{d^{2m}}{dx^{2m}} \left( x^{2r} \right) = \frac{(2r)!}{(2r-2m)!} x^{2r-2m}
$$
Si $m > r$, esta derivada es cero, por lo que $c_{2m} = 0$. Para $m \leq r$, la integral se convierte en una integral gaussiana estándar:
$$
\int_{-\infty}^{\infty} x^{2(r-m)} e^{-x^2} \, dx = \frac{(2(r-m))!}{2^{2(r-m)} (r-m)!} \sqrt{\pi} = \frac{(2r-2m)!}{2^{2r-2m} (r-m)!} \sqrt{\pi}
$$
Sustituyendo este resultado en la expresión para $c_{2m}$:
$$
c_{2m} = \frac{1}{\sqrt{\pi} \, 2^{2m} \, (2m)!} \cdot \frac{(2r)!}{(2r-2m)!} \cdot \frac{(2r-2m)!}{2^{2r-2m} (r-m)!} \sqrt{\pi}
$$
Simplificando los factoriales y las potencias de 2:
$$
c_{2m} = \frac{(2r)!}{2^{2r} \, (2m)! \, (r-m)!}
$$
Cambiando el índice de sumatoria de $m$ a $n$ para coincidir con la notación estándar del enunciado, la expansión de $x^{2r}$ es:
$$
x^{2r} = \frac{(2r)!}{2^{2r}} \sum_{n=0}^{r} \frac{H_{2n}(x)}{(2n)! \, (r-n)!}
$$
Este resultado coincide exactamente con el problema propuesto en la página 340 del texto de Alonso Sepúlveda.

\vspace{0.3cm}
\noindent\textbf{Caso 2: Expansión de $g(x) = x^{2r+1}$} \\
De manera análoga, los coeficientes son:
$$
c_n = \frac{1}{\sqrt{\pi} \, 2^n \, n!} \int_{-\infty}^{\infty} e^{-x^2} x^{2r+1} H_n(x) \, dx
$$
Dado que $x^{2r+1}$ es una función impar, el integrando es par solo cuando $n$ es impar. Por lo tanto, $c_{2m} = 0$. 

Para $n$ impar, hacemos $n = 2m+1$. Usando la fórmula de Rodrigues:
$$
H_{2m+1}(x) = (-1)^{2m+1} e^{x^2} \frac{d^{2m+1}}{dx^{2m+1}} \left( e^{-x^2} \right) = - e^{x^2} \frac{d^{2m+1}}{dx^{2m+1}} \left( e^{-x^2} \right)
$$
Sustituyendo en la integral para $c_{2m+1}$:
$$
c_{2m+1} = \frac{1}{\sqrt{\pi} \, 2^{2m+1} \, (2m+1)!} \int_{-\infty}^{\infty} x^{2r+1} \left[ - \frac{d^{2m+1}}{dx^{2m+1}} \left( e^{-x^2} \right) \right] \, dx
$$
Integrando por partes $2m+1$ veces, los términos de frontera se anulan y obtenemos un factor de $(-1)^{2m+1}$:
$$
c_{2m+1} = \frac{(-1) \cdot (-1)^{2m+1}}{\sqrt{\pi} \, 2^{2m+1} \, (2m+1)!} \int_{-\infty}^{\infty} e^{-x^2} \frac{d^{2m+1}}{dx^{2m+1}} \left( x^{2r+1} \right) \, dx
$$
Dado que $(-1)^{2m+2} = 1$, el signo global es positivo. La derivada $(2m+1)$-ésima de $x^{2r+1}$ es:
$$
\frac{d^{2m+1}}{dx^{2m+1}} \left( x^{2r+1} \right) = \frac{(2r+1)!}{(2r-2m)!} x^{2r-2m}
$$
Nuevamente, si $m > r$, la derivada es cero. Para $m \leq r$, evaluamos la integral gaussiana:
$$
\int_{-\infty}^{\infty} x^{2(r-m)} e^{-x^2} \, dx = \frac{(2r-2m)!}{2^{2r-2m} (r-m)!} \sqrt{\pi}
$$
Sustituyendo en $c_{2m+1}$:
$$
c_{2m+1} = \frac{1}{\sqrt{\pi} \, 2^{2m+1} \, (2m+1)!} \cdot \frac{(2r+1)!}{(2r-2m)!} \cdot \frac{(2r-2m)!}{2^{2r-2m} (r-m)!} \sqrt{\pi}
$$
Simplificando:
$$
c_{2m+1} = \frac{(2r+1)!}{2^{2r+1} \, (2m+1)! \, (r-m)!}
$$
Cambiando el índice de sumatoria de $m$ a $n$, la expansión de $x^{2r+1}$ es:
$$
x^{2r+1} = \frac{(2r+1)!}{2^{2r+1}} \sum_{n=0}^{r} \frac{H_{2n+1}(x)}{(2n+1)! \, (r-n)!}
$$

\vspace{0.3cm}
\noindent\textbf{Discusión y relevancia:} \\
Estas expansiones demuestran cómo cualquier monomio puede expresarse como una combinación lineal finita de polinomios de Hermite. Esto es una consecuencia directa de que $H_n(x)$ es un polinomio de grado $n$ con coeficiente principal $2^n$, lo que garantiza que el conjunto $\{H_n(x)\}$ forma una base completa para el espacio de polinomios. Este tipo de desarrollos es fundamental en mecánica cuántica para calcular elementos de matriz de operadores de posición en la base del oscilador armónico, tal como se discute en la Sección 8.5.2 (p. 341) de Alonso Sepúlveda.

\end{document}
